% ===========================================================================
%  PART 3 --- Deletion, the worst case simulation, and the Fibonacci payoff
%  Flow: deletion rule -> gentle example -> cascade -> Fibonacci tie-back
%        -> rotation bound -> applications -> closing
%
%  Every number is simulated ground truth:
%    worst-case deletion rotations = floor(h/2) exactly, h = 2..14
%    worst-case insertion rotations = 1 at every height tested
%    h=7 tree: deleting the leaf of in-order rank 33 (pre-order f34) fires
%      3 rotations, at ranks 32, 29, 21 (pre f31, f23, f2), depths 3 -> 2 -> 1
%    per-layer subtree sizes: 54 33 20 12 7 4 2 1, and the rotation sites are
%      the T_2, T_4, T_6 layers -- every second one
%  The cascade geometry is regenerated by verify/cascade_gen.py, which asserts
%  all of the above before it writes tikz/cascade.tex. Re-run it after any
%  change to the tree, and never hand-edit the generated file.
% ===========================================================================

% Generated geometry for the deletion cascade below. Loaded HERE rather than
% from avl.tex on purpose, so this part is self-contained: avl.tex needs NO
% change for this to work. These are plain \def's, so loading them inside the
% document body is fine.
%% GENERATED by verify/cascade_gen.py -- do not edit by hand.
%% Worst-case deletion cascade on the minimal height-7 AVL tree.
%% Delete leaf of in-order rank 33 (pre-order f34). Three rotations fire,
% at ranks 32, 29, 21 (pre f31, f23, f2) at depths 3, 2, 1.
%% x of every node is its ORIGINAL in-order rank, held fixed for the whole
%% cascade, so nodes only ever move VERTICALLY -- the same rule Part 2 states.
%% Rank 33 is absent from every state after the first: the node is really gone.

%% --- initial: 54 nodes, height 7
\def\cascNZero{1/7,2/6,3/5,4/6,5/4,6/6,7/5,8/3,9/6,10/5,11/4,12/5,13/2,14/6,15/5,16/4,17/5,18/3,19/5,20/4,21/1,22/6,23/5,24/4,25/5,26/3,27/5,28/4,29/2,30/5,31/4,32/3,33/4,34/0,35/6,36/5,37/4,38/5,39/3,40/5,41/4,42/2,43/5,44/4,45/3,46/4,47/1,48/5,49/4,50/3,51/4,52/2,53/4,54/3}
\def\cascEZero{34/21,21/13,13/8,8/5,5/3,3/2,2/1,3/4,5/7,7/6,8/11,11/10,10/9,11/12,13/18,18/16,16/15,15/14,16/17,18/20,20/19,21/29,29/26,26/24,24/23,23/22,24/25,26/28,28/27,29/32,32/31,31/30,32/33,34/47,47/42,42/39,39/37,37/36,36/35,37/38,39/41,41/40,42/45,45/44,44/43,45/46,47/52,52/50,50/49,49/48,50/51,52/54,54/53}
\def\cascHZero{7}

%% --- before rot 1: 53 nodes, height 7, |bf|=2 at rank 32 (+2)
\def\cascNOne{1/7,2/6,3/5,4/6,5/4,6/6,7/5,8/3,9/6,10/5,11/4,12/5,13/2,14/6,15/5,16/4,17/5,18/3,19/5,20/4,21/1,22/6,23/5,24/4,25/5,26/3,27/5,28/4,29/2,30/5,31/4,32/3,34/0,35/6,36/5,37/4,38/5,39/3,40/5,41/4,42/2,43/5,44/4,45/3,46/4,47/1,48/5,49/4,50/3,51/4,52/2,53/4,54/3}
\def\cascEOne{34/21,21/13,13/8,8/5,5/3,3/2,2/1,3/4,5/7,7/6,8/11,11/10,10/9,11/12,13/18,18/16,16/15,15/14,16/17,18/20,20/19,21/29,29/26,26/24,24/23,23/22,24/25,26/28,28/27,29/32,32/31,31/30,34/47,47/42,42/39,39/37,37/36,36/35,37/38,39/41,41/40,42/45,45/44,44/43,45/46,47/52,52/50,50/49,49/48,50/51,52/54,54/53}
\def\cascHOne{7}
\def\cascHotOne{32}
\def\cascBfOne{+2}

%% --- before rot 2: 53 nodes, height 7, |bf|=2 at rank 29 (+2)
\def\cascNTwo{1/7,2/6,3/5,4/6,5/4,6/6,7/5,8/3,9/6,10/5,11/4,12/5,13/2,14/6,15/5,16/4,17/5,18/3,19/5,20/4,21/1,22/6,23/5,24/4,25/5,26/3,27/5,28/4,29/2,30/4,31/3,32/4,34/0,35/6,36/5,37/4,38/5,39/3,40/5,41/4,42/2,43/5,44/4,45/3,46/4,47/1,48/5,49/4,50/3,51/4,52/2,53/4,54/3}
\def\cascETwo{34/21,21/13,13/8,8/5,5/3,3/2,2/1,3/4,5/7,7/6,8/11,11/10,10/9,11/12,13/18,18/16,16/15,15/14,16/17,18/20,20/19,21/29,29/26,26/24,24/23,23/22,24/25,26/28,28/27,29/31,31/30,31/32,34/47,47/42,42/39,39/37,37/36,36/35,37/38,39/41,41/40,42/45,45/44,44/43,45/46,47/52,52/50,50/49,49/48,50/51,52/54,54/53}
\def\cascHTwo{7}
\def\cascHotTwo{29}
\def\cascBfTwo{+2}

%% --- before rot 3: 53 nodes, height 7, |bf|=2 at rank 21 (+2)
\def\cascNThree{1/7,2/6,3/5,4/6,5/4,6/6,7/5,8/3,9/6,10/5,11/4,12/5,13/2,14/6,15/5,16/4,17/5,18/3,19/5,20/4,21/1,22/5,23/4,24/3,25/4,26/2,27/5,28/4,29/3,30/5,31/4,32/5,34/0,35/6,36/5,37/4,38/5,39/3,40/5,41/4,42/2,43/5,44/4,45/3,46/4,47/1,48/5,49/4,50/3,51/4,52/2,53/4,54/3}
\def\cascEThree{34/21,21/13,13/8,8/5,5/3,3/2,2/1,3/4,5/7,7/6,8/11,11/10,10/9,11/12,13/18,18/16,16/15,15/14,16/17,18/20,20/19,21/26,26/24,24/23,23/22,24/25,26/29,29/28,28/27,29/31,31/30,31/32,34/47,47/42,42/39,39/37,37/36,36/35,37/38,39/41,41/40,42/45,45/44,44/43,45/46,47/52,52/50,50/49,49/48,50/51,52/54,54/53}
\def\cascHThree{7}
\def\cascHotThree{21}
\def\cascBfThree{+2}

%% --- final: 53 nodes, height 6
\def\cascNFour{1/6,2/5,3/4,4/5,5/3,6/5,7/4,8/2,9/5,10/4,11/3,12/4,13/1,14/6,15/5,16/4,17/5,18/3,19/5,20/4,21/2,22/6,23/5,24/4,25/5,26/3,27/6,28/5,29/4,30/6,31/5,32/6,34/0,35/6,36/5,37/4,38/5,39/3,40/5,41/4,42/2,43/5,44/4,45/3,46/4,47/1,48/5,49/4,50/3,51/4,52/2,53/4,54/3}
\def\cascEFour{34/13,13/8,8/5,5/3,3/2,2/1,3/4,5/7,7/6,8/11,11/10,10/9,11/12,13/21,21/18,18/16,16/15,15/14,16/17,18/20,20/19,21/26,26/24,24/23,23/22,24/25,26/29,29/28,28/27,29/31,31/30,31/32,34/47,47/42,42/39,39/37,37/36,36/35,37/38,39/41,41/40,42/45,45/44,44/43,45/46,47/52,52/50,50/49,49/48,50/51,52/54,54/53}
\def\cascHFour{6}

\def\casctarget{33}

\avldivider{3}{Deletion}{The expensive one.}{avlpthree}{avlpthreeA}{avlpthreeB}{avlpthreeC}

\setavlpart{3}{Part 3 --- Deletion}{avlpthree}

\begin{frame}[label=p3]{Deletion: the rule}
  \vspace{0.5em}
  \begin{center}
  \begin{tikzpicture}[x=1cm,y=1cm,
      stepbox/.style={draw=avlpthree!45!avlground,fill=avlpthree!8!avlground,
                   line width=1.5\pxl,rounded corners=6pt,
                   minimum width=3.6cm,minimum height=1.15cm,
                   font=\avlfSec,text=avlink},
      bad/.style ={draw=avlalert!45!avlground,fill=avlalert!8!avlground,
                   line width=1.5\pxl,rounded corners=6pt,
                   minimum width=3.6cm,minimum height=1.15cm,
                   font=\avlfSec,text=avlink}]
    \node[stepbox] (s1) at (-4.5,0)
      {\begin{minipage}{3.2cm}\centering Remove like a normal BST\end{minipage}};
    \node[stepbox] (s2) at (0,0)
      {\begin{minipage}{3.2cm}\centering Walk back up to the root\end{minipage}};
    \node[bad] (s3) at (4.5,0)
      {\begin{minipage}{3.2cm}\centering Rotate at \strong{every} bad node\end{minipage}};
    \draw[avlarrow] (s1) -- (s2);
    \draw[avlarrow] (s2) -- (s3);
  \end{tikzpicture}
  \end{center}

  \vspace{1.0em}
  \begin{center}
    {\avlfLead\color{avlbody}The first two steps are the same as insertion.}\\[0.7em]
    \begin{avlpanelh}[0.62\textwidth]{avlalert}
      {\avlfLead\viol{\wbold The third one is not.}}
    \end{avlpanelh}\\[0.8em]
    {\avlfSec\color{avlmuted}Insertion stops after one single or double repair.
     Deletion may need repairs higher up.}
  \end{center}

  \note{The first two steps look identical to insertion. The difference is the
  last word. Insertion fixes one node and stops. Deletion may have to fix every
  node on the way up.}
\end{frame}

\begin{frame}[label=p3]{Deletion: same rule, new twist}
  \vspace{-0.4em}
  \begin{center}
  \begin{tikzpicture}[x=1cm,y=1cm]
    \path (-5.6,-2.0) rectangle (5.6,1.05);
    \only<1-2>{
      \node[avlannot] at (-3.2,0.65) {delete 10};
      \draw[avledge] (-4.3,-0.8) -- (-3.2,0) -- (-2.1,-0.8);
      \node[avlkey] at (-3.2,0) {20};
      \node[avlkeyalert] at (-4.3,-0.8) {10};
      \node[avlkey] at (-2.1,-0.8) {30};
    }
    \only<2>{
      \draw[avlarrow] (-0.6,-0.45) -- (0.6,-0.45);
      \ghostnode{1.8}{-0.8}{10}
      \draw[avlghostline] (2.9,0) -- (1.8,-0.8);
      \draw[avledge] (2.9,0) -- (4.0,-0.8);
      \node[avlkeyok] at (2.9,0) {20};
      \node[avlkeyok] at (4.0,-0.8) {30};
      \node[avlannotok] at (2.9,0.65) {$\mathrm{bf}(20)=-1$};
      \node[avlannotok] at (2.9,-1.65) {still legal; no rotation};
    }
    \only<3-6>{
      \node[avlannot,anchor=west,align=left,text width=1.8cm] at (-5.3,0.65) {a taller tree\\delete 30};
      \only<3>{
        \draw[avledge] (-2.1,0) -- (-1.0,-0.8);
        \node[avlkeyalert] at (-1.0,-0.8) {30};
      }
      \only<4-6>{
        \ghostnode{-1.0}{-0.8}{30}
        \draw[avlghostline] (-2.1,0) -- (-1.0,-0.8);
      }
      \only<5-6>{\rotright{-2.1,0}{0.9}}
      \draw[avledge] (-2.1,0) -- (-3.2,-0.8) -- (-4.3,-1.6);
      \only<3>{\node[avlkey] at (-2.1,0) {20};}
      \only<4-6>{\node[avlkeyalert] at (-2.1,0) {20};}
      \node[avlkey] at (-3.2,-0.8) {10};
      \node[avlkeygold] at (-4.3,-1.6) {5};
      \only<5-6>{\pivotpin[avlalert]{-2.1,0}}
      \only<4-6>{
        \node[avlannotalert,anchor=east] at (-2.6,-0.1) {$\mathrm{bf}=+2$};
      }
    }
    \only<6>{
      \draw[avlarrow] (-0.4,-0.45) -- (0.8,-0.45);
      \ghostnode{4.0}{0}{20}
      \ghostpath{4.0,-0.3}{4.0,-0.5}
      \ghostnode{2.9}{-0.8}{10}
      \ghostpath{2.9,-0.5}{2.9,-0.3}
      \ghostnode{1.8}{-1.6}{5}
      \ghostpath{1.8,-1.3}{1.8,-1.1}
      \draw[draw=avlbalance!70,line width=1.6\pxl]
        (1.8,-0.8) -- (2.9,0) -- (4.0,-0.8);
      \node[avlkeyok] at (2.9,0) {10};
      \node[avlkeyok] at (1.8,-0.8) {5};
      \node[avlkeyok] at (4.0,-0.8) {20};
      \node[avlannotok] at (2.9,0.65) {one right rotation};
    }
  \end{tikzpicture}

  \begin{minipage}[t][74\pxl][t]{\textwidth}
    \centering\avlfSec\color{avlmuted}
    \only<1>{Start with 20 above 10 and 30.
      Every node is balanced.\par}
    \only<2>{Delete 10. 20 just loses a child --- $\mathrm{bf}(20)=-1$,
      still legal. Nothing to fix.\par}
    \only<3>{A fresh, taller tree. $\mathrm{bf}(20)=+1$, still legal.\\
      This time, delete 30.\par}
    \only<4>{30 is gone. Now $\mathrm{bf}(20)=+2$.\\
      The same alarm from Part 2 goes off.\par}
    \only<5>{Pin 20 and turn clockwise.\\
      One right rotation, just as in Part 2.\par}
    \only<6>{10 becomes the root. The tree is balanced again.\\
      One rotation --- so far, no different from insertion.\par}
  \end{minipage}
  \begin{avlpanelh}[0.94\textwidth]{avlfib}
    \avlfSec\color{avlbody}
    The rule reads the same as insertion.\\
    The difference shows up somewhere else.\par
  \end{avlpanelh}
  \end{center}
  \note{Six clicks, with the before and after separated.
  First, show the balanced three-node tree. Second, delete 10 and reveal the
  result on the right. The balance factor of 20 becomes minus one; no rotation.
  Third, show the fresh taller tree before deleting anything: 20 has children
  10 and 30, and 10 has left child 5. Fourth, delete 30; the balance factor of
  20 becomes plus two. Fifth, introduce the pin and clockwise arc at 20.
  Sixth, reveal the repaired tree on the right: 10 above 5 and 20, with ghosts
  at the old positions. The taller tree is a new example, not a continuation
  of the first deletion.}
\end{frame}

\begin{frame}[label=p3]{From one rotation to a cascade}
  \begin{center}
    {\avlfStmt\wbold\color{avlink}When does a rotation\\
     propagate toward the root?\par}

    \vspace{1.4em}
    {\avlfBody\color{avlmuted}We fixed one node. Is the parent still balanced?\par}
  \end{center}
  \note{Pause here and let the question land. The small example needed only
  one rotation. What could make a repair continue upward instead of stopping?
  Invite the audience to think about what changes for the parent, then show
  the worst-case tree and let the cascade reveal the answer.}
\end{frame}

% ------------------------------------------------------------ WORST CASE SIM
%  The cascade, run for real. Every overlay below is a DIFFERENT tree: the node
%  is actually removed, and each rotation actually restructures the picture.
%  All five states come from verify/cascade_gen.py driving a real AVL
%  implementation; the geometry is read in from tikz/cascade.tex.
%
%  Layout rule (the same one Part 2 states): a node's x is its ORIGINAL
%  in-order rank and never changes. Only depth changes. So a rotation moves
%  nodes straight up or down, the deleted node leaves a permanent gap at its
%  own x, and the height drop at the end is a visible change in the picture
%  rather than a claim in the caption.
\newcommand{\cascdx}{0.196}
\newcommand{\cascdy}{0.50}
\newcommand{\cascrad}{1.5pt}

% \casctree{<nodes>}{<edges>} -- edges behind, discs on top
\newcommand{\casctree}[2]{%
  \foreach \r/\d in #1 {\coordinate (c-\r) at (\r*\cascdx,-\d*\cascdy);}%
  \foreach \a/\b in #2 {\draw[avlgoldedge] (c-\a) -- (c-\b);}%
  \foreach \r/\d in #1 {\fill[avlfib] (c-\r) circle (\cascrad);}}

% The node that is gone. It used to be marked with a dashed ring held at the
% x-slot it would never occupy again; that ring read as a stray artefact once
% the tree had restructured around it, so the slot is now simply left empty.
\newcommand{\cascgap}{}

% \cascbreak{<rank>}{<bf>} -- the node that is now illegal
\newcommand{\cascbreak}[2]{%
  \draw[draw=avlalert, line width=1.2pt] (c-#1) circle (4.2pt);
  % The label goes ABOVE THE ROOT ROW, the one band of this figure that is
  % guaranteed empty at every beat, and a hairline leader ties it back to the
  % ring. Anywhere inside the canopy its filled box lands on some edge and
  % chops it -- which is what it was doing next to the node.
  \draw[avlalert!50, line width=0.5pt]
    ($(c-#1)+(0,4.6pt)$) -- ({#1*\cascdx},0.10);
  \node[text=avlalert, font=\scriptsize\bfseries, anchor=south,
        fill=avlground, inner sep=1.8pt, rounded corners=1pt]
    at ({#1*\cascdx},0.10) {$\mathrm{bf}=#2$};}

% height ruler down the left edge; #1 = height, #2 = colour
\newcommand{\cascruler}[2]{%
  \begin{scope}[xshift=-0.55cm]
    % mirror: the brace must open toward the tree it is measuring
    \draw[decorate,decoration={brace,amplitude=4pt,mirror},#2,line width=0.8pt]
      (0,0.10) -- (0,{-#1*\cascdy-0.10});
    \node[color=#2,font=\scriptsize,anchor=east,rotate=90]
      at (-0.26,{-#1*\cascdy/2}) {height #1};
  \end{scope}}

\begin{frame}[label=p3]{Worst case: one leaf out, and it does not stop}
  \vspace{-0.5em}
  \begin{center}
  \begin{tikzpicture}[x=1cm,y=1cm,scale=0.90,every node/.append style={transform shape}]
    % fixed extent: the picture must not resize between clicks
    \path (-1.5,-3.70) rectangle (11.2,0.60);

    \only<1>{\casctree{\cascNZero}{\cascEZero}%
             \draw[draw=avlalert,line width=1.2pt]
               (\casctarget*\cascdx,-4*\cascdy) circle (3.6pt);
             \cascruler{\cascHZero}{avlfib}}
    \only<2-3>{\casctree{\cascNOne}{\cascEOne}\cascgap
             \cascbreak{\cascHotOne}{\cascBfOne}\cascruler{\cascHOne}{avlfib}}
    \only<4>{\casctree{\cascNTwo}{\cascETwo}\cascgap
             \cascbreak{\cascHotTwo}{\cascBfTwo}\cascruler{\cascHTwo}{avlfib}}
    \only<5>{\casctree{\cascNThree}{\cascEThree}\cascgap
             \cascbreak{\cascHotThree}{\cascBfThree}\cascruler{\cascHThree}{avlfib}}
    \only<6->{\casctree{\cascNFour}{\cascEFour}\cascgap
             \cascruler{\cascHFour}{avlpthree}}

    \only<3>{
      \begin{scope}[overlay]
        \setpinscale{0.22}
        \rotright{c-32}{0.42}
        \pivotpin[avlalert]{c-32}
        \node[avlannotalert,anchor=north west,align=left,text width=3.5cm]
          (shorterNote) at (7.5,-2.65)
          {Unlike insertion: one level shorter. The parent can break too.};
        \draw[avlalert!55,line width=1\pxl]
          ($(c-32)+(0.38,-0.05)$) -- (7.35,-2.48) -- (shorterNote.north west);
      \end{scope}
    }
  \end{tikzpicture}
  \end{center}

  % Caption block reserves its own height: six beats, one of them two lines.
  \vspace{-0.7em}
  \begin{center}
  \begin{minipage}[t][86\pxl][t]{\textwidth}\centering\avlfSec
    \only<1>{{\color{avlmuted}54 nodes, height 7, no spare room anywhere.
              Delete the \viol{marked leaf}.}}
    \only<2-3>{{\color{avlmuted}Gone. Its grandparent is now
              \viol{$\mathrm{bf}=+2$} --- illegal.
              {\color{avlpthree}\wmed Rotate.}}}
    \only<4>{{\color{avlmuted}Fixed --- but that subtree got \emph{shorter},
              so its own parent is now \viol{$+2$}.
              {\color{avlpthree}\wmed Rotate again.}}}
    \only<5>{{\color{avlmuted}And again. This one is a child of the root.
              {\color{avlpthree}\wmed Rotate a third time.}}}
    \only<6->{{\avlfLead\viol{\wbold 3 rotations}\quad\quietc{$\cdot$}\quad
               {\color{avlbody}54 $\rightarrow$ 53 nodes}
               \quad\quietc{$\cdot$}\quad
               {\color{avlbody}height \textbf{7} $\rightarrow$}
               \color{avlpthree}\textbf{6}}\\[0.35em]
              {\color{avlmuted}The height only fell when the cascade reached
               the top --- and one delete paid for all of it.}}
  \end{minipage}
  \end{center}

  \vspace{-0.1em}
  {\avlfAnnot\color{avlmuted}\centering
   Note: this cascade is left-heavy ($\mathrm{bf}=+2$); right-heavy is its mirror.\\
   Height loss can propagate after any Part 2 repair: LL, RR, LR, or RL.\par}

  \note{Keep all five tree states. The extra beat pauses at the first rotation
  trigger so you can point to its pin and arc. This rotation shortens the
  subtree by one --- that is the difference from insertion. A shorter subtree
  can make the parent illegal too. Then continue the original cascade:
  fix the first site, fix its parent, and fix the child of the root.
  The whole tree drops from height seven to six after three rotations.
  This example propagates height loss; not every deletion rotation does.}
\end{frame}

\newcommand{\avldeletionlayers}{%
  \begin{tikzpicture}[x=1cm,y=1cm,
      avldot/.style={circle,draw=avlfib,fill=avlfib,inner sep=0pt,
                     minimum size=5.8\pxl},
      avledge/.style={draw=avlfib!75,line width=1.3\pxl}]
    \setavlscale{0.42}{0.33}
    \resetavl
    \drawfib{7}{0}{0}

    % The root's tag would sit under the title rule if it went left like the
    % rest, so it is the one that leans right instead.
    \node[avlannotfib,anchor=west]
      at ($(f1)+(0.22,0.10)$) {$T_{7}$\,:\,\textbf{54}};
    \foreach \i/\k/\N in {2/6/33, 3/5/20, 4/4/12,
                          5/3/7, 6/2/4, 7/1/2, 8/0/1} {
      \draw[avlfib!45,line width=1\pxl]
        ($(f\i)+(-0.14,0)$) -- ($(f\i)+(-0.30,0)$);
      \node[avlannotfib,anchor=east]
        at ($(f\i)+(-0.32,0.01)$) {$T_{\k}$\,:\,\textbf{\N}};
    }
  \end{tikzpicture}
}

\begin{frame}[t,label=p3]{A Fibonacci tree at every level}
  \vspace{-0.2em}
  \begin{center}
    \avldeletionlayers

    \vspace{0.4em}
    {\avlfSec\color{avlmuted}Each subtree is a smaller \strong{Fibonacci tree}.\\
     It uses the fewest nodes possible for its height.\par}
    \vspace{0.5em}
    {\avlfSec\color{avlfib}$54,\ 33,\ 20,\ 12,\ 7,\ 4,\ 2,\ 1$
     \quad\quietc{---}\quad $N(7)$ down to $N(0)$.\par}
  \end{center}
  \note{The cascade tree is a Fibonacci tree: a minimal AVL tree at every
  level. Each subtree combines the smallest trees of the two preceding
  heights. Its node count follows N(h) = 1 + N(h-1) + N(h-2), so the counts
  are Fibonacci numbers minus one. Read down the left edge: 54, 33, 20, 12,
  7, 4, 2, 1. That is N(7) down to N(0). Next, connect that minimum node
  count to the height loss that drives the cascade.}
\end{frame}

\begin{frame}[t,label=p3]{Why the Fibonacci tree cascades}
  \vspace{-0.2em}
  \begin{center}
    \avldeletionlayers

    \vspace{0.3em}
    {\avlfSec\color{avlmuted}Fibonacci numbers in disguise:
     \fib{$54 = F(10)-1$}.\par}
    \vspace{0.3em}
    {\avlfSec\color{avlmuted}Height 7 needs at least 54 nodes.\\
     Delete one leaf: 53 nodes cannot keep that height.\par}
    \vspace{0.5em}
    \begin{avlpanelh}[0.94\textwidth]{avlpthree}
      \avlfSec\color{avlpthree}No slack anywhere. That's why the fix cascades.\par
    \end{avlpanelh}
  \end{center}
  \note{Same diagram, same thought. The subtree counts are Fibonacci numbers
  minus one: N(h) equals F(h+3) minus one. In particular, 54 is F(10) minus one.
  This tree has no slack. Removing a leaf must reduce its height after repair:
  a height-seven AVL tree needs at least 54 nodes, and only 53 remain.
  Our chosen leaf forces that repair to cascade. Not every leaf must cause
  three rotations; this is the deliberately chosen worst-case deletion.}
\end{frame}

\begin{frame}[label=p3]{Insertion costs 1. Deletion costs $\lfloor h/2 \rfloor$.}
  % Fixed two-column layout: the plot is drawn at its true size (640 px wide)
  % and is never scaled to fit, so its tick and axis labels stay legible.
  % No grid -- the ticks carry the scale, and a guide line under the series
  % labels is the one thing that kept colliding with them.
  \vspace{0.2em}
  \begin{columns}[c,onlytextwidth]
    \begin{column}{7.9cm}
      \begin{tikzpicture}
        \begin{axis}[
          width=7.3cm, height=4.45cm, axis lines=left, scale only axis,
          xlabel={height of the tree}, ylabel={rotations, worst case},
          xlabel style={font=\avlfAnnot,color=avlfaint},
          ylabel style={font=\avlfAnnot,color=avlfaint},
          tick label style={font=\avlfAnnot,color=avlfaint},
          axis line style={color=avlfaint,line width=1\pxl},
          tick style={color=avlfaint,line width=1\pxl},
          xtick={2,4,6,8,10,12,14}, ytick={0,2,4,6,8},
          ymin=0, ymax=8, xmin=2, xmax=14,
        ]
          % the 13 measured points, verify/ ground truth
          \addplot[color=avlalert,line width=2.4\pxl,mark=*,mark size=2\pxl,
                   mark options={fill=avlalert,draw=avlalert}]
            coordinates {(2,1)(3,1)(4,2)(5,2)(6,3)(7,3)(8,4)(9,4)(10,5)(11,5)
                         (12,6)(13,6)(14,7)};
          \addplot[color=avlfib,line width=2.4\pxl,dash pattern=on 8\pxl off 6\pxl]
            coordinates {(2,1)(14,1)};
          % Series named in place, upper left -- no legend box, no swatch of a
          % colour the curve does not already use.
          \node[anchor=west,font=\avlfAnnot\wmed,text=avlalert]
            at (axis cs:2.55,7.3) {\textcolor{avlalert}{\rule[0.6ex]{20\pxl}{2.4\pxl}}\ \ deletion --- $\lfloor h/2 \rfloor$};
          \node[anchor=west,font=\avlfAnnot\wmed,text=avlfib]
            at (axis cs:2.55,6.4) {\textcolor{avlfib}{\rule[0.6ex]{8\pxl}{2.4\pxl}\hskip4\pxl\rule[0.6ex]{8\pxl}{2.4\pxl}}\ \ insertion --- always 1};
        \end{axis}
      \end{tikzpicture}
    \end{column}
    \begin{column}{4.55cm}
      \eyebrow{One million nodes}\\[0.5em]
      {\color{avlrule}\hrule height 1\pxl}\vskip 8\pxl
      \avlstatrow[avlfib]{Insertion}{1}
      \avlstatrow[avlalert]{Deletion}{13}
      \vskip 0.6em
      {\avlfSec\color{avlmuted}Measured on a real AVL tree, every height
       from 2 to 14.}
    \end{column}
  \end{columns}

  \note{Insertion is flat at one, forever. Deletion grows with the height.
  At one million nodes the tree is 27 tall, so a single delete can cost 13
  rotations where an insert costs one.}
\end{frame}

\begin{frame}[label=p3]{Deletion: walk back to the root}
  \begin{center}
    {\avlfStmt\wbold\boldmath\color{avlink}BST deletion: $O(\log n)$.\par}

    \vspace{1.2em}
    {\avlfBody\color{avlmuted}After removing the node, walk upward\\
     from its parent toward the root.\par}

    \vspace{1.2em}
    \begin{avlpanelh}[0.78\textwidth]{avlpthree}
      \avlfLead\color{avlpthree}At most $h=O(\log n)$ ancestors are visited.\par
    \end{avlpanelh}

    \vspace{0.9em}
    {\avlfSec\color{avlmuted}Update the height. Check the balance. Move up.\par}
  \end{center}
  \note{First perform the usual BST deletion. The AVL height bound makes the
  search and successor or predecessor search logarithmic. Start retracing
  from the parent of the node physically removed, which may differ from the
  key initially requested when that key has two children. The path contains
  at most h ancestors, where h is the original tree height. With stored
  heights or balance factors, checking and updating each node costs O(1),
  apart from any rotations.}
\end{frame}

\begin{frame}[label=p3]{Deletion: at most two rotations per level}
  \begin{center}
    {\avlfBody\color{avlmuted}Each unbalanced node needs at most two rotations.\par}

    \vspace{1.0em}
    \begin{minipage}{0.88\textwidth}
      \begin{columns}[T,onlytextwidth]
        \begin{column}{0.47\textwidth}
          \begin{avlpanelh}[\linewidth]{avlfib}
            {\avlfLead\wbold\color{avlfib}LL / RR\par}
            \vspace{0.4em}
            {\avlfBody\color{avlbody}1 rotation\par}
          \end{avlpanelh}
        \end{column}
        \begin{column}{0.47\textwidth}
          \begin{avlpanelh}[\linewidth]{avlfib}
            {\avlfLead\wbold\color{avlfib}LR / RL\par}
            \vspace{0.4em}
            {\avlfBody\color{avlbody}2 rotations\par}
          \end{avlpanelh}
        \end{column}
      \end{columns}
    \end{minipage}

    \vspace{1.1em}
    {\avlfSec\color{avlmuted}At most $h$ ancestors, at most two rotations each:\par}
    \vspace{0.5em}
    {\avlfStmt\color{avlpthree}$\text{Max rotations}\le 2h=O(\log n)$\par}
  \end{center}
  \note{Count primitive rotations, so a double repair counts as two.
  LL and RR need one; LR and RL need two. Deletion may trigger repairs at
  multiple ancestors, but each visited level needs at most one single or
  double repair. Multiplying at most h ancestors by at most two rotations
  gives the safe upper bound 2h. It need not be tight, and many deletions
  require fewer rotations or none.}
\end{frame}

\begin{frame}[label=p3]{AVL deletion stays logarithmic}
  \begin{center}
    {\avlfBody\color{avlmuted}Each rotation takes $O(1)$. Add up the work:\par}

    \vspace{0.8em}
    {\avlfLead\color{avlbody}$\displaystyle
      \underbrace{O(\log n)}_{\text{BST deletion}}
      +\underbrace{O(h)}_{\text{upward walk}}
      +\underbrace{O(h)}_{\text{rotations}}$\par}

    \vspace{1.0em}
    {\avlfStmt\color{avlpthree}$\boxed{\text{AVL Deletion}=O(\log n)}$\par}

    \vspace{1.1em}
    \begin{avlpanelh}[0.96\textwidth]{avlpthree}
      \avlfSec\color{avlbody}Unlike insertion, deletion may rebalance at multiple levels.\\
      Even if repairs reach the root, the total stays $O(\log n)$.\par
    \end{avlpanelh}
  \end{center}
  \note{BST deletion costs O(log n). Visiting at most h ancestors costs O(h).
  At most 2h constant-time rotations also cost O(h), not O(h squared).
  Because h is O(log n), the sum is O(log n). The contrast with insertion is
  the number of repair sites: insertion stops after repairing one site,
  while deletion may continue repairing all the way to% c the root. That extra
  work does not change the logarithmic worst-case bound on the full operation.}
\end{frame}

\begin{frame}[label=p3]{Why bother with the worse constant}
  \begin{center}
  \begin{minipage}{0.90\textwidth}
    \avlfSec\color{avlmuted}
    {\fib{\wbold Databases and filesystems}}\par
    In-memory indexes need guaranteed worst-case lookup,
    not just a good average.\par\vspace{0.9em}
    {\fib{\wbold Read-heavy workloads}}\par
    Lookups vastly outnumber updates. The extra deletion work barely matters.\par\vspace{0.9em}
    {\fib{\wbold A hard bound beats a soft average}}\par
    Real-time systems care about the worst case, not the typical case.\par
  \end{minipage}

  \vspace{1.1em}
  {\avlfSec\color{avlbody}An unbalanced tree is fast until it isn't.
   An AVL tree is always fast\\--- that's the whole trade.\par}
  \end{center}
  \note{AVL trees are useful for ordered in-memory data, not a blanket choice
  for every database or filesystem index. OpenZFS uses AVL trees to track
  in-memory buffers; see include/sys/dnode.h, dn\_dbufs, in the OpenZFS source.
  Disk-oriented database indexes often use B-trees instead; see SQLite's
  Architecture of SQLite, B-Tree section. The read-heavy point is a workload
  trade-off, not a claim that AVL always beats other balanced trees.
  Always fast means logarithmically bounded tree work, not constant time or
  a guaranteed wall-clock deadline. Real-time systems also have to bound
  allocation, synchronization, and the cost of comparing keys.}
\end{frame}

% ------------------------------------------------------------------ closing
\avlstatement{A balanced tree is not one that is even.\\[0.35em]
It is one that is never more than \dstrong{one step} away from even.\\[0.35em]
That one step is what keeps every search fast.}{avlpthreeA}{avlpthreeB}{avlpthreeC}